\scanpage{15}
\mspage{7}
voir que $\operatorname{Centr}_{\pi}(\pi')=\{1\}$, i.e. le

\uline{Corollaire} Soit $\pi$ un groupe profini, extension successive de groupes profinis libres. Alors le centralisateur \inserted{$\unsure{Z|\mathcal Z}$ dans $\pi$} de tout sous-groupe ouvert $\pi'$ de $\pi$ est réduit à $\{1\}$.

Par dévissage on est ramené au cas d'un groupe \inserted{profini} \uline{libre}. \inserted{On sait que $\pi'$ est donc libre.} OPS $\pi'$ invariant, et on admet que \unsure{le|son} centre d'un groupe profini libre est réduit à $1$.
\margin{à vérifier}
Donc $\unsure{Z|\mathcal Z}\longrightarrow\pi/\pi'$ est injectif, donc $\unsure{Z|\mathcal Z}$ est fini, et on admet que dans un groupe profini libre, il n'y a pas d'élément d'ordre fini $\neq1$ --- ce qui achève la démonstration.
\margin{à vérifier}

\uline{Scholie} Le fait que $E_{\bar K/K}$ soit \deleted{\unsure{sans}}\inserted{à} \deleted{\unsure{torsion|tors.}}\inserted{centre} trivial peut s'exploiter, en disant que le groupoïde $\Pi_{\deleted{\bar K/}K}$ \inserted{(ou le topos $B_K$)} est connu à équivalence près, \deleted{\unsure{que l'on|question}}\inserted{défini à isom} unique près, quand on connaît le groupe \uline{extérieur} associé à $E_{\bar K/K}$. Les hom
\margin{14}
